% GENERATED FILE. Edit canonical lessons, then run npm run generate:books.
% canonical-source-hash: fnv1a64:604c2999ba108878
% canonical-lessons: TA-C422-mukam, TA-C422-tarc, TA-C422-kutuvai, TA-C422-tolainokki, TA-C422-avil

\chapter{Camping Gear}
\label{ch:ta-camping-gear}
\hlchaptermodality{422}

\begin{chapteropening}
\textbf{By the end of this chapter:} \emph{I can talk about a camp, a torch, a flask and binoculars, and untying a bundle.}

Complete the last lesson of chapter 422: I can talk about a camp, a torch, a flask and binoculars, and untying a bundle.
\end{chapteropening}

\section[\texorpdfstring{\ta{முகாம்} (mukām)}{முகாம் (mukām)}]{\ta{முகாம்} (mukām) — a camp}

\label{lesson:TA-C422-mukam}



\begin{quote}
Before the new one: say the Tamil for to visit, then the Tamil for to see off.
\end{quote}

\subsection*{You'll want to know: \ta{முகாம்}}
\textbf{\ta{முகாம்}} — \emph{mukām} — \textquotedblleft{}a camp\textquotedblright{}.

\begin{quote}
\ta{மாணவர்கள்} \ta{காட்டில்} \ta{ஒரு} \ta{முகாமில்} \ta{தங்கினார்கள்}.
\end{quote}

\begin{quote}
\emph{māṇavarkaḷ kāṭṭil oru mukāmil taṅkiṉārkaḷ.}
\end{quote}

\begin{quote}
The students stayed at a camp in the forest.
\end{quote}

Before a vowel ending the final \ta{ம்} stays: \ta{முகாமில்}, at the camp.

\begin{cousinweb}
\textbf{\ta{முகாம்}} (\emph{mukām}) is usually traced to Arabic \emph{maqām}, \textquotedblleft{}a place, a halting place\textquotedblright{}, which reached Tamil through Hindustani.
\end{cousinweb}

\begin{grammarlens}[title={What you've built}]
One more word for getting around.
\end{grammarlens}

\subsection*{Your turn}
\begin{itemize}
\raggedright
  \item \emph{Say it:} \emph{mukām}
  \item \emph{Say it:} \emph{mukām}, once more
  \item \emph{Say it:} \emph{māṇavarkaḷ kāṭṭil oru mukāmil taṅkiṉārkaḷ.} — The students stayed at a camp in the forest.
  \item \emph{Recall:} say the Tamil for to visit, then the Tamil for to see off, then say \emph{mukām} again
\end{itemize}

\begin{tcolorbox}[breakable,colback=teal!4,colframe=teal!35!black,title={Before you move on}]
What does \ta{முகாம்} mean? (\textquotedblleft{}A camp\textquotedblright{}.) Say it once more.
\end{tcolorbox}
\section[\texorpdfstring{\ta{டார்ச்} (ṭārc)}{டார்ச் (ṭārc)}]{\ta{டார்ச்} (ṭārc) — a torch; a flashlight}

\label{lesson:TA-C422-tarc}



\begin{quote}
Before the new one: say the Tamil for to see off, then the Tamil for a camp.
\end{quote}

\subsection*{You'll want to know: \ta{டார்ச்}}
\textbf{\ta{டார்ச்}} — \emph{ṭārc} — \textquotedblleft{}a torch; a flashlight\textquotedblright{}.

\begin{quote}
\ta{இருட்டில்} \ta{அவன்} \ta{டார்ச்} \ta{எடுத்துக்கொண்டு} \ta{வெளியே} \ta{போனான்}.
\end{quote}

\begin{quote}
\emph{iruṭṭil avaṉ ṭārc eṭuttukkoṇṭu veḷiyē pōṉāṉ.}
\end{quote}

\begin{quote}
In the dark he took a torch and went out.
\end{quote}

A loan with the final stop's puḷḷi: \ta{டார்ச்}.

\begin{cousinweb}
\textbf{\ta{டார்ச்}} (\emph{ṭārc}) is a loan from English \emph{torch}.
\end{cousinweb}

\begin{grammarlens}[title={What you've built}]
One more word for getting around.
\end{grammarlens}

\subsection*{Your turn}
\begin{itemize}
\raggedright
  \item \emph{Say it:} \emph{ṭārc}
  \item \emph{Say it:} \emph{ṭārc}, once more
  \item \emph{Say it:} \emph{iruṭṭil avaṉ ṭārc eṭuttukkoṇṭu veḷiyē pōṉāṉ.} — In the dark he took a torch and went out.
  \item \emph{Recall:} say the Tamil for to see off, then the Tamil for a camp, then say \emph{ṭārc} again
\end{itemize}

\begin{tcolorbox}[breakable,colback=teal!4,colframe=teal!35!black,title={Before you move on}]
What does \ta{டார்ச்} mean? (\textquotedblleft{}A torch; a flashlight\textquotedblright{}.) Say it once more.
\end{tcolorbox}
\section[\texorpdfstring{\ta{குடுவை} (kuṭuvai)}{குடுவை (kuṭuvai)}]{\ta{குடுவை} (kuṭuvai) — a flask; a narrow-mouthed pot}

\label{lesson:TA-C422-kutuvai}



\begin{quote}
Before the new one: say the Tamil for a camp, then the Tamil for a torch.
\end{quote}

\subsection*{You'll want to know: \ta{குடுவை}}
\textbf{\ta{குடுவை}} — \emph{kuṭuvai} — \textquotedblleft{}a flask; a narrow-mouthed pot\textquotedblright{}.

\begin{quote}
\ta{நான்} \ta{குடுவையில்} \ta{சூடான} \ta{தேநீர்} \ta{கொண்டு} \ta{வந்தேன்}.
\end{quote}

\begin{quote}
\emph{nāṉ kuṭuvaiyil cūṭāṉa tēnīr koṇṭu vantēṉ.}
\end{quote}

\begin{quote}
I brought hot tea in a flask.
\end{quote}

The plural is \ta{குடுவைகள்}.

\begin{cousinweb}
\textbf{\ta{குடுவை}} (\emph{kuṭuvai}) is native Dravidian (DEDR 1651), \textquotedblleft{}a vessel with a small narrow mouth\textquotedblright{}.
\end{cousinweb}

\begin{grammarlens}[title={What you've built}]
One more word for getting around.
\end{grammarlens}

\subsection*{Your turn}
\begin{itemize}
\raggedright
  \item \emph{Say it:} \emph{kuṭuvai}
  \item \emph{Say it:} \emph{kuṭuvai}, once more
  \item \emph{Say it:} \emph{nāṉ kuṭuvaiyil cūṭāṉa tēnīr koṇṭu vantēṉ.} — I brought hot tea in a flask.
  \item \emph{Recall:} say the Tamil for a camp, then the Tamil for a torch, then say \emph{kuṭuvai} again
\end{itemize}

\begin{tcolorbox}[breakable,colback=teal!4,colframe=teal!35!black,title={Before you move on}]
What does \ta{குடுவை} mean? (\textquotedblleft{}A flask; a narrow-mouthed pot\textquotedblright{}.) Say it once more.
\end{tcolorbox}
\section[\texorpdfstring{\ta{தொலைநோக்கி} (tolainōkki)}{தொலைநோக்கி (tolainōkki)}]{\ta{தொலைநோக்கி} (tolainōkki) — binoculars; a telescope}

\label{lesson:TA-C422-tolainokki}



\begin{quote}
Before the new one: say the Tamil for a torch, then the Tamil for a flask.
\end{quote}

\subsection*{You'll want to know: \ta{தொலைநோக்கி}}
\textbf{\ta{தொலைநோக்கி}} — \emph{tolainōkki} — \textquotedblleft{}binoculars; a telescope\textquotedblright{}.

\begin{quote}
\ta{தொலைநோக்கியால்} \ta{தூரத்துப்} \ta{பறவைகளைப்} \ta{பார்த்தோம்}.
\end{quote}

\begin{quote}
\emph{tolainōkkiyāl tūrattup paṟavaikaḷaip pārttōm.}
\end{quote}

\begin{quote}
We watched distant birds through binoculars.
\end{quote}

\ta{தொலை}, far, + \ta{நோக்கி}, a looker; \ta{தொலைக்காட்சி} is built the same way.

\begin{cousinweb}
\textbf{\ta{தொலைநோக்கி}} (\emph{tolainōkki}) is a modern coinage built from \ta{தொலை}, \textquotedblleft{}far\textquotedblright{}, and \ta{நோக்கி}, \textquotedblleft{}one that looks\textquotedblright{}, from \ta{நோக்கு} (DEDR 3794).
\end{cousinweb}

\begin{grammarlens}[title={What you've built}]
One more word for getting around.
\end{grammarlens}

\subsection*{Your turn}
\begin{itemize}
\raggedright
  \item \emph{Say it:} \emph{tolainōkki}
  \item \emph{Say it:} \emph{tolainōkki}, once more
  \item \emph{Say it:} \emph{tolainōkkiyāl tūrattup paṟavaikaḷaip pārttōm.} — We watched distant birds through binoculars.
  \item \emph{Recall:} say the Tamil for a torch, then the Tamil for a flask, then say \emph{tolainōkki} again
\end{itemize}

\begin{tcolorbox}[breakable,colback=teal!4,colframe=teal!35!black,title={Before you move on}]
What does \ta{தொலைநோக்கி} mean? (\textquotedblleft{}Binoculars; a telescope\textquotedblright{}.) Say it once more.
\end{tcolorbox}
\section[\texorpdfstring{\ta{அவிழ்} (aviḻ)}{அவிழ் (aviḻ)}]{\ta{அவிழ்} (aviḻ) — to untie; to undo}

\label{lesson:TA-C422-avil}



\begin{quote}
Before the new one: say the Tamil for a flask, then the Tamil for binoculars.
\end{quote}

\subsection*{You'll want to know: \ta{அவிழ்}}
\textbf{\ta{அவிழ்}} — \emph{aviḻ} — \textquotedblleft{}to untie; to undo\textquotedblright{}.

\begin{quote}
\ta{வீட்டுக்கு} \ta{வந்ததும்} \ta{நான்} \ta{மூட்டையை} \ta{அவிழ்த்தேன்}.
\end{quote}

\begin{quote}
\emph{vīṭṭukku vantatum nāṉ mūṭṭaiyai aviḻttēṉ.}
\end{quote}

\begin{quote}
As soon as I got home, I untied the bundle.
\end{quote}

Past (I …): \ta{அவிழ்த்தேன்}, \emph{aviḻttēṉ}.

\ta{அவிழ்த்தேன்} is I untied it; \ta{அவிழ்ந்தது}, it came undone.

\begin{cousinweb}
\textbf{\ta{அவிழ்}} (\emph{aviḻ}) is native Dravidian (DEDR 14), \textquotedblleft{}to become loose; to loosen, to untie, to unpack\textquotedblright{}.
\end{cousinweb}

\begin{grammarlens}[title={What you've built}]
One more word for getting around.
\end{grammarlens}

\subsection*{Your turn}
\begin{itemize}
\raggedright
  \item \emph{Say it:} \emph{aviḻ}
  \item \emph{Say it:} \emph{aviḻ}, once more
  \item \emph{Say it:} \emph{vīṭṭukku vantatum nāṉ mūṭṭaiyai aviḻttēṉ.} — As soon as I got home, I untied the bundle.
  \item \emph{Recall:} say the Tamil for a flask, then the Tamil for binoculars, then say \emph{aviḻ} again
\end{itemize}

\begin{tcolorbox}[breakable,colback=teal!4,colframe=teal!35!black,title={Before you move on}]
What does \ta{அவிழ்} mean? (\textquotedblleft{}To untie; to undo\textquotedblright{}.) Say it once more.
\end{tcolorbox}
